\section*{Hoofdstuk \ref{ch:b2:structure-determination} --- Structuurbepaling: \texorpdfstring{$^{13}$C}{13C}-NMR, MS en IR samen}

\begin{solution}{exo:b2:structure-determination:1}
1,2-Dimethylbenzeen: 4 (\ce{CH3}; C1/C2; C3/C6; C4/C5). 1,3-: 5 (\ce{CH3}; C1/C3; C2; C4/C6; C5).
1,4-: 3 (\ce{CH3}; C1/C4; C2/C3/C5/C6).
\end{solution}

\begin{solution}{exo:b2:structure-determination:2}
\ce{CH3CH(OH)CH2CH3}. DEPT-135: C1 (\ce{CH3}) omhoog, C2 (CH) omhoog, C3 (\ce{CH2}) omlaag, C4 (\ce{CH3}) omhoog.
DEPT-90: alleen C2.
\end{solution}

\begin{solution}{exo:b2:structure-determination:3}
209: ketoncarbonyl; 170: ester- (of zuur)carbonyl; 129: \omterm{def:b2:huckel:aromatic}{aromatisch} koolstofatoom; 64: koolstofatoom gebonden aan een
zuurstofatoom; 14: alkyl-\ce{CH3}.
\end{solution}

\begin{solution}{exo:b2:structure-determination:4}
$(2 \times 8 + 2 - 8)/2 = 5$: een benzeenring (4) en een \ce{C=O} (1). Bijvoorbeeld methylbenzoaat,
\ce{C6H5COOCH3}, of 4-methylbenzoëzuur.
\end{solution}

\begin{solution}{exo:b2:structure-determination:5}
Eén onverzadiging, een keton (\qty{1715}{cm^{-1}}, lijn bij 209, quaternair); één \ce{CH2} en twee
\ce{CH3}, vier verschillende koolstofatomen: butaan-2-on, \ce{CH3COCH2CH3}.
\end{solution}

\begin{solution}{exo:b2:structure-determination:6}
Koolstof: pentaan-2-on toont vijf lijnen, het symmetrische pentaan-3-on drie. Protonen: pentaan-2-on heeft een
singlet van drie protonen (\ce{CH3CO}); pentaan-3-on alleen een kwartet en een triplet. \omterm{def:b2:mass-spec-atomic:spectrum}{Massaspectrum}: pentaan-2-on
heeft zijn \omterm{def:b2:mass-spec-atomic:spectrum}{basispiek} bij 43 (\ce{CH3CO+}) en een McLafferty-ion bij 58; pentaan-3-on heeft zijn \omterm{def:b2:mass-spec-atomic:spectrum}{basispiek} bij 57
(\ce{C2H5CO+}).
\end{solution}

\begin{solution}{exo:b2:structure-determination:7}
\ce{C4H8O2}, 88: het paar kwartet--triplet is een \ce{OCH2CH3} (de \ce{CH2} bij 4.12 zit op zuurstof), het
singlet een \ce{CH3CO}. Ethylethanoaat, \ce{CH3COOCH2CH3}.
\end{solution}

\begin{solution}{exo:b2:structure-determination:8}
1,2-Dichloorbenzeen: drie lijnen (C1/C2, C3/C6, C4/C5). 1,4-Dichloorbenzeen: twee (C1/C4, de vier
CH). DEPT-90 toont twee CH-lijnen voor het eerste, één voor het tweede.
\end{solution}

\begin{solution}{exo:b2:structure-determination:9}
C2 is een stereocentrum. Het ene of het andere proton van C3 door een deuteriumatoom vervangen geeft twee
diastereo-isomeren, geen enantiomeren: de twee protonen zijn diastereotoop, in verschillende omgevingen.
Ze kunnen verschillende verschuivingen hebben en zowel met elkaar als met hun buren koppelen, zodat de
\ce{CH2} een complex multiplet geeft in plaats van een eenvoudig kwintet.
\end{solution}

\begin{solution}{exo:b2:structure-determination:10}
Onverzadigingsgraad 5; een monogesubstitueerde ring (drie CH-lijnen en één C-lijn bij 137.0, vijf \omterm{def:b2:huckel:aromatic}{aromatische}
H); een keton (200.8, C); een ethylgroep (\ce{CH2} 31.8, kwartet bij 2.98 naast de carbonylgroep, \ce{CH3}
8.3). 1-Fenylpropaan-1-on, \ce{C6H5COCH2CH3}. Conjugatie met de ring delokaliseert de $\pi$-elektronen van \ce{C=O}
en verzwakt de binding: de band zakt onder de \qty{1715}{cm^{-1}} van een verzadigd keton.
% ledger: ir:CO-ketone
\end{solution}

\begin{solution}{exo:b2:structure-determination:11}
Pentaanzuur: de zeer brede band van \ce{O-H} van 2500 tot \qty{3300}{cm^{-1}} en een proton rond 11--12
ppm. Methylbutanoaat: een singlet van drie protonen in het gebied van protonen op een koolstofatoom gebonden aan zuurstof
(3.3--4.5). Ethylpropanoaat: twee kwartetten, één op zuurstof en één naast de carbonylgroep.
Propylethanoaat: een singlet van drie protonen naast de carbonylgroep (2.0--2.4) en een triplet op zuurstof.
% ledger: ir:OH-acid, nmrr:acid, nmrr:alcohol, nmrr:ketone
\end{solution}

\begin{solution}{exo:b2:structure-determination:12}
Beide lijnen verdwijnen bij \omterm{def:b2:structure-determination:dept}{DEPT}, dus beide zijn quaternair, en \omterm{def:b2:structure-determination:dept}{DEPT} alleen kan ze niet onderscheiden. De verschuivingen
kunnen dat wel: estercarbonylen liggen rond 167--171 en \omterm{def:b2:huckel:aromatic}{aromatische} koolstofatomen tussen 128 en 137 in de grafiek van dit
hoofdstuk. Een ringkoolstofatoom bij 166.5 en een carbonyl bij 130.6 zouden allebei ver buiten hun klassen liggen; de
juiste toekenning is de omgekeerde.
\end{solution}

\begin{solution}{pb:b2:structure-determination:1}
\textbf{1.} $3.1/31.7 \approx 9.8~\%$; $9.8/1.11 \approx 9$ koolstofatomen.
\textbf{2.} Even massa: geen stikstof (of twee stikstofatomen).
\textbf{3.} $150 - 108 = 42$: \ce{H10O2} past (10 + 32): \ce{C9H10O2}.
\textbf{4.} Tien koolstofatomen zouden $M+1 \approx 11~\%$ van $M$ geven, niet 9.8~\%; en het IR-spectrum toont een
carbonylgroep, die naast de band van \ce{C-O} een tweede zuurstofatoom vraagt.
\textbf{5.} $(2 \times 9 + 2 - 10)/2 = 5$.
\textbf{6.} $108 + 10 \times 1.007825 + 2 \times 15.994915 \approx 150.0681$.
\textbf{7.} Een strekvibratie van \ce{C=O} op de positie van een verzadigde ester (rond \qty{1735}{cm^{-1}}).
\textbf{8.} De strekvibratie van \ce{C-O} van de ester.
\textbf{9.} \omterm{def:b2:huckel:aromatic}{Aromatische} (of alkeen-)bindingen \ce{C-H}.
\textbf{10.} Een groep \ce{O-H}: geen alcohol, geen carbonzuur.
\textbf{11.} Nee: een carbonylgroep die met een ring geconjugeerd is, absorbeert bij een lager golfgetal; hier ligt ze waar
verzadigde esters liggen, dus een niet-conjugerende groep scheidt ze van de ring.
\textbf{12.} 170.70: ester-\ce{C=O}; 136.14: ringkoolstofatoom dat de substituent draagt; 128.56 en 128.24:
\omterm{def:b2:huckel:aromatic}{aromatische} CH; 66.24: \ce{CH2} gebonden aan zuurstof; 20.82: \ce{CH3}.
\textbf{13.} 66.24, de \ce{CH2}.
\textbf{14.} 170.70 en 136.14, de twee koolstofatomen zonder waterstof.
\textbf{15.} 128.56 en 128.24.
\textbf{16.} Vier: het koolstofatoom dat de substituent draagt, de twee ortho-, de twee meta- en de para-CH. Er worden drie
soorten CH verwacht; twee ervan hebben verschuivingen die te dicht bij elkaar liggen om bij dit veld gescheiden te worden, en delen een lijn.
\textbf{17.} Niet noodzakelijk: hoogten zijn geen aantallen (\cref{prop:b2:structure-determination:intensities}),
al kan de lijn hier best twee overlappende soorten CH bevatten.
\textbf{18.} 7.33: de vijf \omterm{def:b2:huckel:aromatic}{aromatische} protonen; 5.09: de \ce{OCH2}; 2.06: de \ce{CH3CO}.
\textbf{19.} Hun naburige atomen (het ringkoolstofatoom en het zuurstofatoom voor de \ce{CH2}, het carbonylkoolstofatoom
voor de \ce{CH3}) dragen geen waterstof.
\textbf{20.} Ze is gebonden aan het elektronegatieve zuurstofatoom van de ester en aan de ring: beide ontschermen haar.
\textbf{21.} \ce{C6H5CH2OCOCH3}, benzylethanoaat.
\textbf{22.} In het isomeer is de \ce{CH2} niet aan zuurstof gebonden, dus ze kan niet bij 5.09 verschijnen, en de
methylgroep, op zuurstof, zou in het gebied 3.3--4.5 verschijnen, niet bij 2.06 naast een carbonylgroep.
% ledger: nmrr:alcohol, nmrr:ketone
\textbf{23.} 108: verlies van 42 (keteen, \ce{CH2=C=O}) door omlegging, het radicaalkation van
\ce{C7H8O}; 91: \ce{C7H7+}, het benzylkation; 43: \ce{CH3CO+}.
\textbf{24.} Zeven soorten koolstofatomen (\ce{C=O}, het gesubstitueerde ringkoolstofatoom, ortho-, meta- en para-CH,
\ce{CH2}, \ce{CH3}): \textbf{zeven lijnen, waarvan één (de \ce{CH2}) omlaag wijst} bij DEPT-135.
\end{solution}
